\paragraph{What the diagnostics show.} Three findings recur across our
experiments. First, the supracompetitive outcomes of Q-learning informed
traders in the standard Kyle market do not rest on punishment: rivals do not
react to deviations even with perfect monitoring, and deviating pays. Second,
the same learning procedure under-trades with no rival at all, by an amount
that grows with $\sqrt\alpha$ and disappears under counterfactual updates.
Third, detectability, not learning sophistication, decides whether punishment
is even feasible: in the standard Kyle market a deviation is a fraction of a
noise standard deviation (Proposition~\ref{prop:snr}), whereas with many
information-insensitive investors it is hundreds of standard deviations.

\paragraph{Policy.} If supracompetitive trading is a learning bias, the
interventions that target collusion through information, such as delayed or
coarser order-flow disclosure, will not work, and our reduced-transparency
treatment confirms this. What does work is algorithm design:
counterfactual (full-information) updating removes the bias. This mirrors the
conclusion of \citet{asker2022,asker2024} and \citet{abada2023} for pricing
algorithms and extends it to informed trading. It also suggests a practical
audit for supervisors: before treating a supracompetitive outcome as
collusion, retrain the algorithm myopically and check whether the outcome
survives. A behaviour that a myopic learner reproduces cannot be a punishment
scheme.

\paragraph{Limitations.} Our evidence concerns tabular Q-learning and two deep
RL algorithms with standard hyperparameters; more sophisticated learners, for
instance ones that model their rivals, may find punishment strategies that
these do not. Convergence of constant-step Q-learning is not complete in the
strict sense of \citet{calvano2020}: greedy strategies still change on path by
about one grid step over the last $10^5$ periods, although outcomes are stable
across seeds, schedules and horizons. Our market maker learns its price impact
from exponentially weighted moments rather than being fixed at the
equilibrium value, which is realistic but differs from some prior designs; our
\citet{dou2025} replication follows their calibration, grid and state, but not
every implementation detail (for example their state-dependent exploration
schedule and training horizons of up to $5\times10^{10}$ periods). Finally,
the stylised market abstracts from inventory dynamics, multiple assets and
long-lived information, where the scope for monitoring and punishment may
differ.

\paragraph{Conclusion.} Low trading intensity by learning algorithms is
easy to produce and easy to misread. Telling collusion from learning bias
requires diagnostics that can fail: a deviation test, a placebo that cannot
collude, and a check of whether the outcome moves with the learning rule. In
the canonical market for informed trading, these diagnostics point to pruning,
not punishment.
